\chapter{Results: The Paper's Other Tasks}\label{ch:results-other}

\section{NL-to-Scala}\label{sec:results-scala}

All 152 cells are scored: 19 Instruct models $\times$ 4 test sets $\times$ 2 prompts. Scoring finished on 2026-11-16
at 06:04. The Scala adapters were trained inside the NL-to-SQL runs of \cref{ch:results-sql,ch:results-ext}, on the
Scala side of the same validated (and, for repair-TS, recovered) pairs.

{\footnotesize
\rowcolors{3}{rowshade}{white}
\setlength{\tabcolsep}{4.2pt}
\begin{longtable}{@{}lccrrrrrrrrrrr@{}}
\caption{NL-to-Scala, all 19 Instruct models: execution accuracy per test set and prompt (repo-faithful, following the repo's \texttt{nss\_evaluator}), with the corrected zero-shot mean.}\label{tab:gen-scala}\\
\toprule
\rowcolor{headshade}Model & S & R & \multicolumn{2}{c}{WikiSQL} & \multicolumn{2}{c}{WikiTQ} & \multicolumn{2}{c}{Spider} & \multicolumn{2}{c}{BIRD} & M8 & ZS4 & ZS4 corr.\\\rowcolor{headshade} & & & z & f & z & f & z & f & z & f & & &\\
\midrule
\endfirsthead
\multicolumn{14}{l}{\textit{\tablename~\thetable{} (continued)}}\\
\toprule
\rowcolor{headshade}Model & S & R & \multicolumn{2}{c}{WikiSQL} & \multicolumn{2}{c}{WikiTQ} & \multicolumn{2}{c}{Spider} & \multicolumn{2}{c}{BIRD} & M8 & ZS4 & ZS4 corr.\\\rowcolor{headshade} & & & z & f & z & f & z & f & z & f & & &\\
\midrule
\endhead
\midrule
\multicolumn{14}{r}{\textit{continued on next page}}\\
\endfoot
\bottomrule
\endlastfoot
Vanilla & -- & -- & 0.3410 & 0.6720 & 0.0420 & 0.1290 & 0.3005 & 0.3907 & 0.0909 & 0.1625 & 0.2661 & 0.1936 & 0.2146\\
\midrule
TS & 1 & 1 & 0.6940 & 0.6720 & 0.1930 & 0.2250 & 0.4674 & 0.5292 & 0.2287 & 0.2204 & 0.4037 & 0.3958 & 0.4201\\
TS & 1 & 2 & 0.7160 & 0.6790 & 0.2900 & 0.3070 & 0.5735 & 0.5860 & 0.2507 & 0.2590 & 0.4577 & 0.4576 & 0.4806\\
TS & 1 & 3 & 0.7190 & 0.6870 & 0.3290 & 0.2950 & 0.5743 & 0.5743 & 0.3030 & 0.2617 & 0.4679 & 0.4813 & 0.5050\\
TS & 2 & 1 & 0.7320 & 0.7220 & 0.1720 & 0.2060 & 0.4349 & 0.4800 & 0.2011 & 0.1956 & 0.3929 & 0.3850 & 0.4085\\
TS & 2 & 2 & 0.6850 & 0.6790 & 0.2900 & 0.2780 & 0.5601 & 0.5317 & 0.3003 & 0.2645 & 0.4486 & 0.4589 & 0.4840\\
TS & 2 & 3 & 0.6290 & 0.6330 & 0.3260 & 0.3030 & 0.5776 & 0.5710 & 0.2975 & 0.2755 & 0.4516 & 0.4575 & 0.4804\\
TS & 3 & 1 & 0.7150 & 0.6530 & 0.1580 & 0.1820 & 0.4407 & 0.4733 & 0.1956 & 0.1928 & 0.3763 & 0.3773 & 0.4014\\
TS & 3 & 2 & 0.6880 & 0.6900 & 0.2060 & 0.2220 & 0.5092 & 0.5217 & 0.2231 & 0.2149 & 0.4094 & 0.4066 & 0.4317\\
TS & 3 & 3 & 0.6650 & 0.6690 & 0.2900 & 0.2950 & 0.5776 & 0.5952 & 0.3113 & 0.2893 & 0.4616 & 0.4610 & 0.4834\\
\midrule
Repair-TS & 1 & 1 & 0.6960 & 0.6770 & 0.2320 & 0.2280 & 0.5058 & 0.5317 & 0.2479 & 0.2452 & 0.4204 & 0.4204 & 0.4457\\
Repair-TS & 1 & 2 & 0.7250 & 0.6970 & 0.3390 & 0.3380 & 0.6210 & 0.6311 & 0.2893 & 0.2617 & 0.4878 & 0.4936 & 0.5193\\
Repair-TS & 1 & 3 & 0.6500 & 0.6520 & 0.3480 & 0.3530 & 0.6252 & 0.6202 & 0.2893 & 0.2782 & 0.4770 & 0.4781 & 0.5039\\
Repair-TS & 2 & 1 & 0.7510 & 0.7260 & 0.1990 & 0.1970 & 0.4407 & 0.4833 & 0.2011 & 0.1901 & 0.3985 & 0.3980 & 0.4229\\
Repair-TS & 2 & 2 & 0.7230 & 0.6960 & 0.3210 & 0.2910 & 0.5902 & 0.5626 & 0.2672 & 0.2479 & 0.4624 & 0.4753 & 0.5010\\
Repair-TS & 2 & 3 & 0.6920 & 0.6680 & 0.3420 & 0.3150 & 0.6269 & 0.6060 & 0.2727 & 0.2810 & 0.4754 & 0.4834 & 0.5109\\
Repair-TS & 3 & 1 & 0.6650 & 0.6790 & 0.1800 & 0.2160 & 0.4791 & 0.5109 & 0.2204 & 0.2204 & 0.3964 & 0.3861 & 0.4104\\
Repair-TS & 3 & 2 & 0.7030 & 0.6800 & 0.3080 & 0.3030 & 0.5935 & 0.5776 & 0.3085 & 0.2590 & 0.4666 & 0.4783 & 0.5021\\
Repair-TS & 3 & 3 & 0.6990 & 0.6980 & 0.3210 & 0.3180 & 0.6219 & 0.6085 & 0.2920 & 0.2727 & 0.4789 & 0.4835 & 0.5065\\
\end{longtable}
}
\vspace{-0.6\baselineskip}
{\footnotesize\color{slate}\textit{Note.} Scala adapters were trained inside the NL-to-SQL runs; no new training. Source: \fpath{results/baseline/final_nlscala_instruct_*.json}.\par}
\medskip

{\small
\rowcolors{3}{rowshade}{white}
\setlength{\tabcolsep}{4.2pt}
\begin{longtable}{@{}lccrrrrrrrr@{}}
\caption{NL-to-Scala summary statistics over seeds (repo-faithful).}\label{tab:gen-scala-stats}\\
\toprule
\rowcolor{headshade}Approach & Rd & Seeds & \multicolumn{4}{c}{M8 over seeds} & \multicolumn{4}{c}{ZS4 over seeds}\\\rowcolor{headshade} & & & Mean & SD & Min & Max & Mean & SD & Min & Max\\
\midrule
\endfirsthead
\multicolumn{11}{l}{\textit{\tablename~\thetable{} (continued)}}\\
\toprule
\rowcolor{headshade}Approach & Rd & Seeds & \multicolumn{4}{c}{M8 over seeds} & \multicolumn{4}{c}{ZS4 over seeds}\\\rowcolor{headshade} & & & Mean & SD & Min & Max & Mean & SD & Min & Max\\
\midrule
\endhead
\midrule
\multicolumn{11}{r}{\textit{continued on next page}}\\
\endfoot
\bottomrule
\endlastfoot
Vanilla & -- & 1 & 0.2661 & -- & 0.2661 & 0.2661 & 0.1936 & -- & 0.1936 & 0.1936\\
TS & 1 & 3 & 0.3910 & 0.0138 & 0.3763 & 0.4037 & 0.3860 & 0.0093 & 0.3773 & 0.3958\\
TS & 2 & 3 & 0.4385 & 0.0257 & 0.4094 & 0.4577 & 0.4410 & 0.0298 & 0.4066 & 0.4589\\
TS & 3 & 3 & 0.4603 & 0.0082 & 0.4516 & 0.4679 & 0.4666 & 0.0129 & 0.4575 & 0.4813\\
Repair-TS & 1 & 3 & 0.4051 & 0.0133 & 0.3964 & 0.4204 & 0.4015 & 0.0174 & 0.3861 & 0.4204\\
Repair-TS & 2 & 3 & 0.4722 & 0.0136 & 0.4624 & 0.4878 & 0.4824 & 0.0098 & 0.4753 & 0.4936\\
Repair-TS & 3 & 3 & 0.4771 & 0.0017 & 0.4754 & 0.4789 & 0.4817 & 0.0031 & 0.4781 & 0.4835\\
\end{longtable}
}
\vspace{-0.6\baselineskip}
{\footnotesize\color{slate}\textit{Note.} Recalculated from the per-cell files.\par}
\medskip


\begin{figure}[htbp]
\centering
\begin{tikzpicture}
\begin{axis}[width=0.62\textwidth,height=6.0cm,xmin=0.7,xmax=3.3,xtick={1,2,3},xlabel={Round},
  ylabel={\ZSiv{} (zero-shot mean of 4)},ymin=0.15,ymax=0.52,ymajorgrids,grid style={headshade},
  legend style={font=\scriptsize,at={(0.97,0.03)},anchor=south east,draw=none},tick label style={font=\scriptsize}]
  \addplot[navy,thick,mark=*,error bars/.cd,y dir=both,y explicit] coordinates {\plotScalaTS};
  \addplot[sea,thick,mark=triangle*,error bars/.cd,y dir=both,y explicit] coordinates {\plotScalaRTS};
  \addplot[slate,dashed,domain=0.7:3.3] {\valScalaVanilla};
  \addplot[amber,thick,dotted,domain=0.7:3.3] {0.3750};
  \addplot[amber!60!black,densely dashed,domain=0.7:3.3] {0.1878};
  \legend{TS,repair-TS,our vanilla,paper TS,paper vanilla}
\end{axis}
\end{tikzpicture}
\caption{NL-to-Scala \ZSiv{} by round (mean over seeds, one-SD error bars), against the paper's Llama
figures (horizontal lines; published Table~3, mean of the 4 public sets).}
\label{fig:scala}
\end{figure}

\begin{observation}
\begin{itemize}
\item \textbf{Vanilla} has \ZSiv{} 0.1936 against the paper's 0.1878. It is within 0.018 on WikiSQL
(0.341), WikiTQ (0.042) and BIRD (0.091), and $+0.052$ on Spider (0.3005).
\item \textbf{TS} improves every round with both prompts: \ZSiv{} 0.3860, 0.4410, 0.4666. At round~3 it is 0.0916
above the paper's TS (0.3750). It is higher on WikiTQ ($+0.127$), Spider ($+0.221$) and BIRD ($+0.076$),
and 0.057 lower on WikiSQL.
\item \textbf{Repair-TS} is ahead of TS at every round on \Mviii{}: $+0.0141$, $+0.0337$ and $+0.0167$, all
$p<0.0001$, and all 9 seed-round differences positive (\cref{tab:gen-significance_scala}).
\item On \ZSiv{}, repair-TS peaks by round~2 (0.4824) and is flat at round~3 (0.4817), with a quarter of TS's seed
spread (SD 0.0031 against 0.0129).
\item Its gains are on WikiTQ and Spider. BIRD at round~3 is $-0.019$.
\item \textbf{Few-shot} helps vanilla most (\ZSiv{} 0.1936; few-shot mean of 4 0.3386) and slightly lowers the
fine-tuned models: TS 0.4541 and repair-TS 0.4726 few-shot, against 0.4666 and 0.4817 zero-shot.
\end{itemize}
\end{observation}

\begin{interpretation}
Two results stand out.
\begin{itemize}
\item \textbf{TS on Scala is well above the paper's figure.} We cannot attribute this: corpus, model
checkpoint and the paper's round choice are all candidate causes.
\item \textbf{Repair-TS helps Scala even on Instruct, where it does not help SQL} (\cref{sec:ext-repair}). On
Instruct, 60\% of recovered pairs (450 of 744) carry a Scala program written by the repairer, against 35\% for SQL
(259). Repair therefore adds more new Scala than new SQL to the training sets. Scala is also the weaker language:
vanilla scores 0.19 against 0.50 for SQL (\ZSiv{}). Both facts are consistent with the result, but neither is tested
as its cause.
\end{itemize}
\end{interpretation}

\section{Error detection}\label{sec:results-ed}

{\footnotesize
\rowcolors{3}{rowshade}{white}
\setlength{\tabcolsep}{4.2pt}
\begin{longtable}{@{}lcrrrrr@{}}
\caption{Error detection, summary: repo-faithful F1 per seed, mean $\pm$ sample SD over seeds, and the mean of ST and RT.}\label{tab:gen-ed-summary}\\
\toprule
\rowcolor{headshade}Model & R & \makecell{ST seeds\\1 / 2 / 3} & ST mean $\pm$ SD & \makecell{RT seeds\\1 / 2 / 3} & RT mean $\pm$ SD & Avg\\
\midrule
\endfirsthead
\multicolumn{7}{l}{\textit{\tablename~\thetable{} (continued)}}\\
\toprule
\rowcolor{headshade}Model & R & \makecell{ST seeds\\1 / 2 / 3} & ST mean $\pm$ SD & \makecell{RT seeds\\1 / 2 / 3} & RT mean $\pm$ SD & Avg\\
\midrule
\endhead
\midrule
\multicolumn{7}{r}{\textit{continued on next page}}\\
\endfoot
\bottomrule
\endlastfoot
Paper vanilla & -- & -- & 0.071 & -- & 0.108 & 0.090\\
Paper TS & final & -- & 0.136 & -- & 0.161 & 0.148\\
\midrule
Vanilla & -- & -- & 0.0659 & -- & 0.1111 & 0.0885\\
\midrule
TS & 1 & 0.1397 / 0.2211 / 0.1677 & 0.1762 $\pm$ 0.0414 & 0.1989 / 0.2836 / 0.2466 & 0.2430 $\pm$ 0.0425 & 0.2096\\
TS & 2 & 0.1954 / 0.2609 / 0.2600 & 0.2388 $\pm$ 0.0376 & 0.2742 / 0.3009 / 0.3265 & 0.3005 $\pm$ 0.0262 & 0.2697\\
TS & 3 & 0.2655 / 0.2909 / 0.2791 & 0.2785 $\pm$ 0.0127 & 0.3191 / 0.2857 / 0.2889 & 0.2979 $\pm$ 0.0184 & 0.2882\\
\midrule
Repair-ED & 1 & 0.1545 / 0.2429 / 0.2034 & 0.2003 $\pm$ 0.0443 & 0.2353 / 0.3303 / 0.2553 & 0.2736 $\pm$ 0.0501 & 0.2369\\
Repair-ED & 2 & 0.3077 / 0.2430 / 0.1916 & 0.2474 $\pm$ 0.0582 & 0.3366 / 0.3333 / 0.2439 & 0.3046 $\pm$ 0.0526 & 0.2760\\
Repair-ED & 3 & 0.2459 / 0.2558 / 0.2742 & 0.2586 $\pm$ 0.0144 & 0.2752 / 0.3768 / 0.3673 & 0.3398 $\pm$ 0.0561 & 0.2992\\
\midrule
Best-of-2 & 1 & 0.1311 / 0.2156 / 0.2286 & 0.1918 $\pm$ 0.0529 & 0.1834 / 0.3443 / 0.2837 & 0.2705 $\pm$ 0.0813 & 0.2311\\
\midrule
Same-data & 1 & 0.1561 / 0.2404 / 0.1792 & 0.1919 $\pm$ 0.0436 & 0.2280 / 0.2258 / 0.2532 & 0.2357 $\pm$ 0.0152 & 0.2138\\
\end{longtable}
}
\vspace{-0.6\baselineskip}
{\footnotesize\color{slate}\textit{Note.} Paper figures: published Table~3, Llama-3.1-8B. Source: \fpath{results/baseline/train_ed_*.json}.\par}
\medskip

{\footnotesize
\rowcolors{3}{rowshade}{white}
\setlength{\tabcolsep}{4.2pt}
\begin{longtable}{@{}lcclrrrrrrr@{}}
\caption{Error detection on AutoTest (Llama-3.1-8B-Instruct, repo zero-shot prompt). Precision is averaged over columns with a prediction, recall over columns with errors, and F1 is taken from the two averages (the repo's \texttt{ed\_evaluator}); the AutoTest protocol ignores predictions of debatable values.}\label{tab:gen-ed}\\
\toprule
\rowcolor{headshade}Model & S & R & Bench & Columns & \makecell{With\\errors} & Predicted & Precision & Recall & F1 (repo) & \makecell{F1\\(AutoTest)}\\
\midrule
\endfirsthead
\multicolumn{11}{l}{\textit{\tablename~\thetable{} (continued)}}\\
\toprule
\rowcolor{headshade}Model & S & R & Bench & Columns & \makecell{With\\errors} & Predicted & Precision & Recall & F1 (repo) & \makecell{F1\\(AutoTest)}\\
\midrule
\endhead
\midrule
\multicolumn{11}{r}{\textit{continued on next page}}\\
\endfoot
\bottomrule
\endlastfoot
Vanilla & -- & -- & ST & 1126 & 42 & 504 & 0.0357 & 0.4286 & 0.0659 & 0.0694\\
Vanilla & -- & -- & RT & 1081 & 36 & 342 & 0.0614 & 0.5833 & 0.1111 & 0.1151\\
\midrule
TS & 1 & 1 & ST & 1126 & 42 & 187 & 0.0856 & 0.3810 & 0.1397 & 0.1531\\
TS & 1 & 1 & RT & 1081 & 36 & 145 & 0.1241 & 0.5000 & 0.1989 & 0.2045\\
TS & 1 & 2 & ST & 1126 & 42 & 132 & 0.1288 & 0.4048 & 0.1954 & 0.2179\\
TS & 1 & 2 & RT & 1081 & 36 & 88 & 0.1932 & 0.4722 & 0.2742 & 0.2857\\
TS & 1 & 3 & ST & 1126 & 42 & 71 & 0.2113 & 0.3571 & 0.2655 & 0.2778\\
TS & 1 & 3 & RT & 1081 & 36 & 58 & 0.2586 & 0.4167 & 0.3191 & 0.3261\\
TS & 2 & 1 & ST & 1126 & 42 & 148 & 0.1419 & 0.5000 & 0.2211 & 0.2386\\
TS & 2 & 1 & RT & 1081 & 36 & 98 & 0.1939 & 0.5278 & 0.2836 & 0.2923\\
TS & 2 & 2 & ST & 1126 & 42 & 96 & 0.1875 & 0.4286 & 0.2609 & 0.2769\\
TS & 2 & 2 & RT & 1081 & 36 & 77 & 0.2208 & 0.4722 & 0.3009 & 0.3091\\
TS & 2 & 3 & ST & 1126 & 42 & 68 & 0.2353 & 0.3810 & 0.2909 & 0.3168\\
TS & 2 & 3 & RT & 1081 & 36 & 55 & 0.2364 & 0.3611 & 0.2857 & 0.3023\\
TS & 3 & 1 & ST & 1126 & 42 & 125 & 0.1120 & 0.3333 & 0.1677 & 0.1795\\
TS & 3 & 1 & RT & 1081 & 36 & 110 & 0.1636 & 0.5000 & 0.2466 & 0.2553\\
TS & 3 & 2 & ST & 1126 & 42 & 58 & 0.2241 & 0.3095 & 0.2600 & 0.2737\\
TS & 3 & 2 & RT & 1081 & 36 & 62 & 0.2581 & 0.4444 & 0.3265 & 0.3333\\
TS & 3 & 3 & ST & 1126 & 42 & 44 & 0.2727 & 0.2857 & 0.2791 & 0.2927\\
TS & 3 & 3 & RT & 1081 & 36 & 54 & 0.2407 & 0.3611 & 0.2889 & 0.2921\\
\midrule
Repair-ED & 1 & 1 & ST & 1126 & 42 & 178 & 0.0955 & 0.4048 & 0.1545 & 0.1700\\
Repair-ED & 1 & 1 & RT & 1081 & 36 & 134 & 0.1493 & 0.5556 & 0.2353 & 0.2454\\
Repair-ED & 1 & 2 & ST & 1126 & 42 & 88 & 0.2273 & 0.4762 & 0.3077 & 0.3509\\
Repair-ED & 1 & 2 & RT & 1081 & 36 & 65 & 0.2615 & 0.4722 & 0.3366 & 0.3542\\
Repair-ED & 1 & 3 & ST & 1126 & 42 & 80 & 0.1875 & 0.3571 & 0.2459 & 0.2609\\
Repair-ED & 1 & 3 & RT & 1081 & 36 & 73 & 0.2055 & 0.4167 & 0.2752 & 0.2885\\
Repair-ED & 2 & 1 & ST & 1126 & 42 & 98 & 0.1735 & 0.4048 & 0.2429 & 0.2576\\
Repair-ED & 2 & 1 & RT & 1081 & 36 & 73 & 0.2466 & 0.5000 & 0.3303 & 0.3429\\
Repair-ED & 2 & 2 & ST & 1126 & 42 & 65 & 0.2000 & 0.3095 & 0.2430 & 0.2549\\
Repair-ED & 2 & 2 & RT & 1081 & 36 & 48 & 0.2917 & 0.3889 & 0.3333 & 0.3373\\
Repair-ED & 2 & 3 & ST & 1126 & 42 & 44 & 0.2500 & 0.2619 & 0.2558 & 0.2651\\
Repair-ED & 2 & 3 & RT & 1081 & 36 & 33 & 0.3939 & 0.3611 & 0.3768 & 0.3824\\
Repair-ED & 3 & 1 & ST & 1126 & 42 & 135 & 0.1333 & 0.4286 & 0.2034 & 0.2195\\
Repair-ED & 3 & 1 & RT & 1081 & 36 & 105 & 0.1714 & 0.5000 & 0.2553 & 0.2628\\
Repair-ED & 3 & 2 & ST & 1126 & 42 & 125 & 0.1280 & 0.3810 & 0.1916 & 0.2051\\
Repair-ED & 3 & 2 & RT & 1081 & 36 & 128 & 0.1562 & 0.5556 & 0.2439 & 0.2500\\
Repair-ED & 3 & 3 & ST & 1126 & 42 & 82 & 0.2073 & 0.4048 & 0.2742 & 0.2881\\
Repair-ED & 3 & 3 & RT & 1081 & 36 & 62 & 0.2903 & 0.5000 & 0.3673 & 0.3830\\
\midrule
Best-of-2 & 1 & 1 & ST & 1126 & 42 & 202 & 0.0792 & 0.3810 & 0.1311 & 0.1404\\
Best-of-2 & 1 & 1 & RT & 1081 & 36 & 193 & 0.1088 & 0.5833 & 0.1834 & 0.1867\\
Best-of-2 & 2 & 1 & ST & 1126 & 42 & 125 & 0.1440 & 0.4286 & 0.2156 & 0.2338\\
Best-of-2 & 2 & 1 & RT & 1081 & 36 & 86 & 0.2442 & 0.5833 & 0.3443 & 0.3621\\
Best-of-2 & 3 & 1 & ST & 1126 & 42 & 133 & 0.1504 & 0.4762 & 0.2286 & 0.2439\\
Best-of-2 & 3 & 1 & RT & 1081 & 36 & 105 & 0.1905 & 0.5556 & 0.2837 & 0.2920\\
\midrule
Same-data & 1 & 1 & ST & 1126 & 42 & 227 & 0.0925 & 0.5000 & 0.1561 & 0.1707\\
Same-data & 1 & 1 & RT & 1081 & 36 & 157 & 0.1401 & 0.6111 & 0.2280 & 0.2378\\
Same-data & 2 & 1 & ST & 1126 & 42 & 141 & 0.1560 & 0.5238 & 0.2404 & 0.2588\\
Same-data & 2 & 1 & RT & 1081 & 36 & 88 & 0.1591 & 0.3889 & 0.2258 & 0.2333\\
Same-data & 3 & 1 & ST & 1126 & 42 & 170 & 0.1118 & 0.4524 & 0.1792 & 0.1939\\
Same-data & 3 & 1 & RT & 1081 & 36 & 122 & 0.1639 & 0.5556 & 0.2532 & 0.2614\\
\end{longtable}
}
\vspace{-0.6\baselineskip}
{\footnotesize\color{slate}\textit{Note.} The controls (best-of-2, same-data) were run for round~1 only. Source: \fpath{results/baseline/train_ed_*.json}, \fpath{final_ed_instruct_vanilla_*.json}.\par}
\medskip

{\small
\rowcolors{3}{rowshade}{white}
\setlength{\tabcolsep}{4.2pt}
\begin{longtable}{@{}lccrrrrrr@{}}
\caption{Error-detection data per round: Algorithm~2 outcomes (kept / candidates; all 5 shuffles must agree), recovered errors for the repair and resample arms, and training examples.}\label{tab:gen-ed-rounds}\\
\toprule
\rowcolor{headshade}Arm & S & R & Typo kept & \makecell{Semantic\\kept} & \makecell{Clean\\kept} & \makecell{Recovered\\(typo + sem.)} & Examples & \makecell{Training\\h}\\
\midrule
\endfirsthead
\multicolumn{9}{l}{\textit{\tablename~\thetable{} (continued)}}\\
\toprule
\rowcolor{headshade}Arm & S & R & Typo kept & \makecell{Semantic\\kept} & \makecell{Clean\\kept} & \makecell{Recovered\\(typo + sem.)} & Examples & \makecell{Training\\h}\\
\midrule
\endhead
\midrule
\multicolumn{9}{r}{\textit{continued on next page}}\\
\endfoot
\bottomrule
\endlastfoot
TS & 1 & 1 & 477/2007 & 663/2569 & 1048/2000 & -- & 2188 & 2.61\\
TS & 1 & 2 & 491/1983 & 1393/2573 & 1653/2000 & -- & 3537 & 0.85\\
TS & 1 & 3 & 541/1993 & 1649/2580 & 1746/2000 & -- & 3936 & 2.69\\
TS & 2 & 1 & 513/2075 & 727/2578 & 1056/2000 & -- & 2296 & 2.56\\
TS & 2 & 2 & 785/2066 & 1582/2574 & 1703/2000 & -- & 4070 & 3.15\\
TS & 2 & 3 & 698/2056 & 1793/2569 & 1786/2000 & -- & 4277 & 3.13\\
TS & 3 & 1 & 475/1993 & 696/2571 & 1085/2000 & -- & 2256 & 1.94\\
TS & 3 & 2 & 423/1991 & 1368/2582 & 1740/2000 & -- & 3531 & 3.06\\
TS & 3 & 3 & 476/2003 & 1585/2577 & 1884/2000 & -- & 3945 & 3.39\\
Repair-ED & 1 & 1 & 477/2007 & 663/2569 & 1048/2000 & 178 + 328 & 2694 & 2.31\\
Repair-ED & 1 & 2 & 480/1983 & 1283/2573 & 1700/2000 & 184 + 307 & 3954 & 3.45\\
Repair-ED & 1 & 3 & 543/1993 & 1433/2580 & 1823/2000 & 201 + 452 & 4452 & 2.55\\
Repair-ED & 2 & 1 & 513/2075 & 727/2578 & 1056/2000 & 175 + 316 & 2787 & 1.49\\
Repair-ED & 2 & 2 & 672/2066 & 1514/2574 & 1728/2000 & 207 + 366 & 4487 & 2.47\\
Repair-ED & 2 & 3 & 553/2056 & 1554/2569 & 1838/2000 & 267 + 385 & 4597 & 2.32\\
Repair-ED & 3 & 1 & 475/1993 & 696/2571 & 1085/2000 & 151 + 367 & 2774 & 1.27\\
Repair-ED & 3 & 2 & 468/1991 & 1596/2582 & 1744/2000 & 253 + 443 & 4504 & 2.21\\
Repair-ED & 3 & 3 & 591/2003 & 1678/2577 & 1827/2000 & 284 + 402 & 4782 & 2.42\\
Best-of-2 & 1 & 1 & 477/2007 & 663/2569 & 1048/2000 & 168 + 278 & 2634 & 1.76\\
Best-of-2 & 2 & 1 & 513/2075 & 727/2578 & 1056/2000 & 199 + 247 & 2742 & 1.49\\
Best-of-2 & 3 & 1 & 475/1993 & 696/2571 & 1085/2000 & 165 + 266 & 2687 & 1.45\\
Same-data & 1 & 1 & -- & -- & -- & 506 & 2694 & 1.53\\
Same-data & 2 & 1 & -- & -- & -- & 491 & 2787 & 1.51\\
Same-data & 3 & 1 & -- & -- & -- & 518 & 2774 & 1.42\\
\end{longtable}
}
\vspace{-0.6\baselineskip}
{\footnotesize\color{slate}\textit{Note.} Clean columns come from the first 1,000 tables of the round, 2 per table. Source: \fpath{runs/train_ed*_instruct_seed<s>/round<r>/summary.json}.\par}
\medskip


\begin{figure}[htbp]
\centering
\begin{tikzpicture}
\begin{groupplot}[group style={group size=2 by 1,horizontal sep=1.5cm},width=0.47\textwidth,height=5.8cm,
  xmin=0.7,xmax=3.3,xtick={1,2,3},xlabel={Round},ymin=0,ymax=0.4,ymajorgrids,grid style={headshade},
  legend style={font=\tiny,at={(0.97,0.03)},anchor=south east,draw=none,fill=white,fill opacity=0.85},
  tick label style={font=\scriptsize},title style={font=\small\bfseries}]
\nextgroupplot[title={ST-bench (F1, repo metric)},ylabel={F1}]
  \addplot[navy,thick,mark=*] coordinates {\plotEDTSSTsa};
  \addplot[navy!60,thick,mark=square*] coordinates {\plotEDTSSTsb};
  \addplot[navy!35,thick,mark=diamond*] coordinates {\plotEDTSSTsc};
  \addplot[coral,thick,mark=triangle*] coordinates {\plotEDREDSTsa};
  \addplot[coral!60,thick,mark=triangle] coordinates {\plotEDREDSTsb};
  \addplot[slate,dashed,domain=0.7:3.3] {0.0659};
  \addplot[amber,dotted,thick,domain=0.7:3.3] {0.136};
  \legend{TS s1,TS s2,TS s3,repair-ED s1,repair-ED s2,our vanilla,paper TS\textsuperscript{\dag}}
\nextgroupplot[title={RT-bench (F1, repo metric)}]
  \addplot[navy,thick,mark=*] coordinates {\plotEDTSRTsa};
  \addplot[navy!60,thick,mark=square*] coordinates {\plotEDTSRTsb};
  \addplot[navy!35,thick,mark=diamond*] coordinates {\plotEDTSRTsc};
  \addplot[coral,thick,mark=triangle*] coordinates {\plotEDREDRTsa};
  \addplot[coral!60,thick,mark=triangle] coordinates {\plotEDREDRTsb};
  \addplot[slate,dashed,domain=0.7:3.3] {0.1111};
  \addplot[amber,dotted,thick,domain=0.7:3.3] {0.161};
\end{groupplot}
\end{tikzpicture}
\caption{Error-detection F1 per seed and round. Repair-ED lines show the rounds finished when the tables were
generated.}
\label{fig:ed}
\end{figure}

\begin{observation}
\textbf{Vanilla} (Instruct, repo zero-shot prompt) has F1 0.0659 (ST) and 0.1111 (RT), against the paper's 0.071 and
0.108: within 0.006. It flags about 40\% of columns, almost all of them clean (precision 0.036 and 0.061),
and finds about half of the real errors (recall 0.43 and 0.58). 36 ST and 46 RT answers had no parseable JSON, which
counts as no prediction.

\textbf{TS}, round~3 means over 3 seeds: ST 0.2785, RT 0.2979. That is about twice the paper's
0.136 and 0.161.
\begin{itemize}
\item ST rises every round on every seed.
\item RT rises every round on seed~1 only. Seeds 2 and 3 fall from round~2 to round~3 (0.3009 $\to$ 0.2857 and 0.3265
$\to$ 0.2889).
\end{itemize}

\textbf{The validator becomes more lenient each round} (\cref{tab:gen-ed-rounds}). For seed~1:
\begin{itemize}
\item semantic errors kept: 663, then 1,393, then 1,649;
\item clean columns kept: 1,048, then 1,653, then 1,746 (of 2,000);
\item training examples: 2,188, then 3,537, then 3,936.
\end{itemize}
\end{observation}

\begin{observation}
\textbf{Repair-ED and its controls} (\cref{tab:gen-ed-summary}; tests in \cref{tab:gen-significance-ed}). Each round
recovered 491--696 rejected errors (\cref{tab:gen-ed-rounds}).
\begin{itemize}
\item \textbf{Round~1:} repair-ED averages 0.2369 against TS's 0.2096, $+0.027$ ($p=0.03$, all three seeds positive).
\item \textbf{Same-data control} (as many extra kept errors from another seed's TS round~1): 0.2138, only $+0.004$ over
TS ($p=0.69$). Repair-ED beats it by $+0.023$ ($p=0.03$; RT $+0.038$, $p=0.009$), so the round-1 gain is not data
volume alone.
\item \textbf{Best-of-2 control} (a second sample without feedback, then the same validation): 0.2311, $+0.022$ over TS
($p=0.12$). Repair-ED is only $+0.006$ above it ($p=0.71$), with seeds split in sign.
\item \textbf{Rounds 2 and 3:} the gain over TS shrinks to $+0.006$ ($p=0.70$) and $+0.011$ ($p=0.58$). At round~3
repair-ED is lower on ST ($-0.020$) and higher on RT ($+0.042$), and seed~1 is below TS.
\end{itemize}
\end{observation}

\begin{interpretation}
Retrying rejected errors helps the first round, but the evidence does not show that the execution feedback is what
helps: a plain second sample does about as well. The effect does not accumulate over rounds, unlike repair-TS on
NL-to-Scala and NL-to-R. The test sets are small (42 and 36 columns with errors), so differences of 0.02--0.04 F1
are within one or two columns, and every interval except round~1's crosses zero.
\end{interpretation}

\begin{interpretation}
Our error-detection TS exceeds the paper's TS figure by a wide margin, from a vanilla that matches it.
Two of our loop choices are not in the paper: vanilla writes the errors, and the previous round's adapter validates.
Either could raise the result, as could our corpus. The round-3 RT decline on two seeds coincides with a validator that
keeps more columns each round. One possible reading is that later validators accept less obvious errors, which makes
the training signal noisier. This is a hypothesis, not tested.
\end{interpretation}

\section{NL-to-R}\label{sec:results-r}

NL-to-R was approved on 2026-11-16 at 08:40. It was built the same day, under the user's instruction not to slow the
running error-detection experiments.

\begin{itemize}
\item \textbf{R sandbox.} Our \fpath{sandbox/Dockerfile.r} is the repository's recipe (r-base by digest, user 65534)
plus Debian's \texttt{r-cran-dplyr}, which every repository NL-to-R prompt loads. Installing it upgraded R from 4.4.2
to 4.6.1 (dplyr 1.2.1). One R program takes about 2.3--2.7~s in the sandbox, against about 40--46~s for Spark.
\item \textbf{Code.} Validation (\fpath{ts/validate/r_agree.py}), the loop (\fpath{ts/r_loop.py}, TS and repair-TS
arms) and scoring (\fpath{ts/eval/r_eval.py}, the repository's \texttt{nr\_evaluator} rules). 9 unit tests; the full
suite passes.
\item \textbf{CPU checks.}
  \begin{itemize}
  \item Hand-written R for three WikiSQL test cases scored as expected in the real sandbox: two correct, one deliberately
  wrong.
  \item A real TS pair with a deliberately wrong R program was correctly rejected (\texttt{disagree} on the first
  sample), and its repair feedback was correctly formed.
  \end{itemize}
\item \textbf{Schedule.} GPU runs use GPU~0 only until error detection finishes, with the user's approval of
2026-11-16 10:05.
\end{itemize}

\textbf{Runs.} TS-R and repair-TS-R, 3 seeds $\times$ 3 rounds each, 3,000 tables per round, finished on
2026-11-17. The smoke test exposed a bug before the full runs: repair feedback kept only the first line of dplyr's
errors (\texttt{Error in filter():}) and dropped the cause on the following lines; it was fixed and unit-tested
(\cref{sec:fail-r-feedback}). Validation kept 321--579 pairs per round (about 11--19\% of questions); repair recovered
84--126 more: 25--31\% on top in round~1 and 16--23\% in rounds 2--3.

{\footnotesize
\rowcolors{3}{rowshade}{white}
\setlength{\tabcolsep}{4.2pt}
\begin{longtable}{@{}lccrrrrrrrrrrr@{}}
\caption{NL-to-R, all 19 Instruct models: execution accuracy per test set and prompt (repo-faithful, following the repo's \texttt{nr\_evaluator}), with the corrected zero-shot mean.}\label{tab:gen-nlr}\\
\toprule
\rowcolor{headshade}Model & S & R & \multicolumn{2}{c}{WikiSQL} & \multicolumn{2}{c}{WikiTQ} & \multicolumn{2}{c}{Spider} & \multicolumn{2}{c}{BIRD} & M8 & ZS4 & ZS4 corr.\\\rowcolor{headshade} & & & z & f & z & f & z & f & z & f & & &\\
\midrule
\endfirsthead
\multicolumn{14}{l}{\textit{\tablename~\thetable{} (continued)}}\\
\toprule
\rowcolor{headshade}Model & S & R & \multicolumn{2}{c}{WikiSQL} & \multicolumn{2}{c}{WikiTQ} & \multicolumn{2}{c}{Spider} & \multicolumn{2}{c}{BIRD} & M8 & ZS4 & ZS4 corr.\\\rowcolor{headshade} & & & z & f & z & f & z & f & z & f & & &\\
\midrule
\endhead
\midrule
\multicolumn{14}{r}{\textit{continued on next page}}\\
\endfoot
\bottomrule
\endlastfoot
Vanilla & -- & -- & 0.3260 & 0.3830 & 0.1500 & 0.2220 & 0.5993 & 0.6152 & 0.2314 & 0.2782 & 0.3506 & 0.3267 & 0.3357\\
\midrule
TS-R & 1 & 1 & 0.4550 & 0.4650 & 0.3020 & 0.2860 & 0.7037 & 0.6861 & 0.3581 & 0.3223 & 0.4473 & 0.4547 & 0.4701\\
TS-R & 1 & 2 & 0.4780 & 0.5320 & 0.3080 & 0.3230 & 0.7312 & 0.6962 & 0.3774 & 0.3444 & 0.4738 & 0.4737 & 0.4919\\
TS-R & 1 & 3 & 0.5850 & 0.6210 & 0.3310 & 0.3580 & 0.7329 & 0.7220 & 0.3609 & 0.3471 & 0.5072 & 0.5024 & 0.5202\\
TS-R & 2 & 1 & 0.3980 & 0.4310 & 0.2850 & 0.2790 & 0.6970 & 0.6586 & 0.3416 & 0.3388 & 0.4286 & 0.4304 & 0.4488\\
TS-R & 2 & 2 & 0.4580 & 0.4950 & 0.3050 & 0.3140 & 0.7129 & 0.6870 & 0.3636 & 0.3609 & 0.4621 & 0.4599 & 0.4778\\
TS-R & 2 & 3 & 0.4370 & 0.4760 & 0.3060 & 0.3120 & 0.7053 & 0.6753 & 0.3719 & 0.3416 & 0.4531 & 0.4551 & 0.4730\\
TS-R & 3 & 1 & 0.3910 & 0.4110 & 0.2580 & 0.2760 & 0.6811 & 0.6694 & 0.3361 & 0.3278 & 0.4188 & 0.4166 & 0.4292\\
TS-R & 3 & 2 & 0.4320 & 0.4860 & 0.2800 & 0.2890 & 0.6861 & 0.6728 & 0.3609 & 0.3416 & 0.4435 & 0.4398 & 0.4573\\
TS-R & 3 & 3 & 0.4950 & 0.5290 & 0.3070 & 0.3250 & 0.7087 & 0.7062 & 0.3719 & 0.3609 & 0.4755 & 0.4707 & 0.4882\\
\midrule
Repair-TS-R & 1 & 1 & 0.4570 & 0.4800 & 0.3120 & 0.3020 & 0.7112 & 0.6945 & 0.3361 & 0.3085 & 0.4502 & 0.4541 & 0.4685\\
Repair-TS-R & 1 & 2 & 0.5290 & 0.5280 & 0.3290 & 0.3360 & 0.7321 & 0.7129 & 0.3636 & 0.3361 & 0.4833 & 0.4884 & 0.5068\\
Repair-TS-R & 1 & 3 & 0.6100 & 0.6340 & 0.3480 & 0.3620 & 0.7346 & 0.7321 & 0.3802 & 0.3664 & 0.5209 & 0.5182 & 0.5346\\
Repair-TS-R & 2 & 1 & 0.4090 & 0.4310 & 0.2950 & 0.2890 & 0.7137 & 0.6878 & 0.3581 & 0.3416 & 0.4406 & 0.4439 & 0.4610\\
Repair-TS-R & 2 & 2 & 0.5220 & 0.5540 & 0.3310 & 0.3450 & 0.7154 & 0.6953 & 0.3691 & 0.3471 & 0.4849 & 0.4844 & 0.5032\\
Repair-TS-R & 2 & 3 & 0.6240 & 0.6100 & 0.3190 & 0.3170 & 0.6903 & 0.6920 & 0.3388 & 0.3251 & 0.4895 & 0.4930 & 0.5100\\
Repair-TS-R & 3 & 1 & 0.4150 & 0.4710 & 0.2860 & 0.2940 & 0.7012 & 0.6694 & 0.3278 & 0.3306 & 0.4369 & 0.4325 & 0.4480\\
Repair-TS-R & 3 & 2 & 0.4810 & 0.5420 & 0.3090 & 0.3080 & 0.7162 & 0.7012 & 0.3691 & 0.3471 & 0.4717 & 0.4688 & 0.4877\\
Repair-TS-R & 3 & 3 & 0.5730 & 0.5820 & 0.3330 & 0.3440 & 0.7229 & 0.6987 & 0.3774 & 0.3581 & 0.4986 & 0.5016 & 0.5215\\
\end{longtable}
}
\vspace{-0.6\baselineskip}
{\footnotesize\color{slate}\textit{Note.} R adapters trained by \fpath{ts/r_loop.py} (SQL--R validation). Source: \fpath{results/baseline/final_nlr_instruct_*.json}.\par}
\medskip

{\small
\rowcolors{3}{rowshade}{white}
\setlength{\tabcolsep}{4.2pt}
\begin{longtable}{@{}lccrrrrrrrr@{}}
\caption{NL-to-R summary statistics over seeds (repo-faithful).}\label{tab:gen-nlr-stats}\\
\toprule
\rowcolor{headshade}Approach & Rd & Seeds & \multicolumn{4}{c}{M8 over seeds} & \multicolumn{4}{c}{ZS4 over seeds}\\\rowcolor{headshade} & & & Mean & SD & Min & Max & Mean & SD & Min & Max\\
\midrule
\endfirsthead
\multicolumn{11}{l}{\textit{\tablename~\thetable{} (continued)}}\\
\toprule
\rowcolor{headshade}Approach & Rd & Seeds & \multicolumn{4}{c}{M8 over seeds} & \multicolumn{4}{c}{ZS4 over seeds}\\\rowcolor{headshade} & & & Mean & SD & Min & Max & Mean & SD & Min & Max\\
\midrule
\endhead
\midrule
\multicolumn{11}{r}{\textit{continued on next page}}\\
\endfoot
\bottomrule
\endlastfoot
Vanilla & -- & 1 & 0.3506 & -- & 0.3506 & 0.3506 & 0.3267 & -- & 0.3267 & 0.3267\\
TS-R & 1 & 3 & 0.4316 & 0.0145 & 0.4188 & 0.4473 & 0.4339 & 0.0193 & 0.4166 & 0.4547\\
TS-R & 2 & 3 & 0.4598 & 0.0152 & 0.4435 & 0.4738 & 0.4578 & 0.0170 & 0.4398 & 0.4737\\
TS-R & 3 & 3 & 0.4786 & 0.0272 & 0.4531 & 0.5072 & 0.4760 & 0.0242 & 0.4551 & 0.5024\\
Repair-TS-R & 1 & 3 & 0.4426 & 0.0068 & 0.4369 & 0.4502 & 0.4435 & 0.0108 & 0.4325 & 0.4541\\
Repair-TS-R & 2 & 3 & 0.4800 & 0.0072 & 0.4717 & 0.4849 & 0.4805 & 0.0103 & 0.4688 & 0.4884\\
Repair-TS-R & 3 & 3 & 0.5030 & 0.0161 & 0.4895 & 0.5209 & 0.5043 & 0.0128 & 0.4930 & 0.5182\\
\end{longtable}
}
\vspace{-0.6\baselineskip}
{\footnotesize\color{slate}\textit{Note.} Recalculated from the per-cell files.\par}
\medskip


\begin{observation}
\begin{itemize}
\item \textbf{Vanilla} has \ZSiv{} 0.3267 (few-shot 0.3746). Most failures are code that does not run, typically column
names containing spaces left unquoted in \texttt{dplyr} calls. The paper's Table~3 gives 0.331, 0.138, 0.503 and
0.225 on WikiSQL, WikiTQ, Spider and BIRD (mean of 4: 0.2993; its 0.260 average includes Text2Analysis). Ours is within
0.012 on WikiSQL, WikiTQ and BIRD, and $+0.096$ on Spider, the same Spider offset seen for NL-to-SQL.
\item \textbf{TS-R} improves every round: \ZSiv{} 0.4339, 0.4578, 0.4760. At round~3 it is $+0.128$ over vanilla on
\Mviii{} ($p<0.0001$, every seed), and round~3 beats round~1 by $+0.047$ ($p<0.0001$, every seed). The largest gains
are on WikiTQ (0.150 $\to$ 0.315) and WikiSQL (0.326 $\to$ 0.506).
\item \textbf{Against the paper's TS} (0.409, 0.257, 0.536, 0.261; mean of 4 0.3658), our round-3 TS-R is higher on all
four sets: WikiSQL $+0.097$, WikiTQ $+0.058$, Spider $+0.180$, BIRD $+0.107$, and $+0.110$ on the mean of 4. The paper's
own gain over vanilla is $+0.067$ on the same four sets; ours is $+0.149$.
\item \textbf{Repair-TS-R} beats TS-R at every round on every seed: \Mviii{} $+0.0110$, $+0.0202$, $+0.0244$, all
$p<0.0001$ (\cref{tab:gen-significance_r}). Its round-3 \ZSiv{} is 0.5043 $\pm$ 0.0128, the best NL-to-R model, with
half TS-R's seed SD. The gain comes mostly from WikiSQL ($+0.097$ zero-shot at round~3) and WikiTQ; Spider and BIRD
are flat.
\end{itemize}
\end{observation}

\begin{interpretation}
NL-to-R behaves like NL-to-Scala, not like Instruct NL-to-SQL: the trained language's first-pass data is weak (only
11\% of round-1 questions pass validation), rounds help, and recovered data helps more each round. This fits the
conditional claim of \cref{sec:novelty}. No same-data control was run for R, so the share of the gain due to the extra
16--31\% of pairs, rather than to their repaired content, is not measured.
\end{interpretation}
